% Generated draft: requires mathematical and source review.
\documentclass[11pt]{article}
\usepackage[margin=25mm]{geometry}
\usepackage{fontspec}
\setmainfont{Latin Modern Roman}
\usepackage{amsmath,amssymb,mathtools,unicode-math}
\setmathfont{Latin Modern Math}
\usepackage{enumitem,needspace,xurl,hyperref}
\setlength{\parindent}{0pt}
\setlength{\parskip}{5pt}
\setlength{\emergencystretch}{3em}
\newcommand{\N}{\mathbb{N}}
\newcommand{\Z}{\mathbb{Z}}
\newcommand{\Q}{\mathbb{Q}}
\newcommand{\R}{\mathbb{R}}
\newcommand{\C}{\mathbb{C}}
\newcommand{\F}{\mathbb{F}}
\newcommand{\A}{\mathbb{A}}
\newcommand{\PP}{\mathbb{P}}
\newcommand{\E}{\mathbb{E}}
\newcommand{\eps}{\varepsilon}
\newcommand{\dd}{\,\mathrm{d}}
\begin{document}
\section*{UnsolvedMath ID: 30000506}
\textbf{Source problem code:} catalogue entry
\subsection{Definitions and mathematical statement}
Let \(S\) be the surface of the intended stretch-line construction and \(\mathcal T(S)\) its Teichm\"uller space. In the usual finite-type hyperbolic setting, points are marked hyperbolic structures, and a stretch line \(t\mapsto X_t\) is the directed deformation determined by an initial structure and a suitable complete geodesic lamination, with exponential stretch factor under the customary parametrization. The source does not specify the surface type, punctures or boundary, lamination hypotheses, or normalization of \(t\).

Thurston's compactification has boundary \(\mathbb P\mathcal{ML}(S)\), the projective classes of nonzero measured laminations in that usual setting. Let \(\mathcal C(S)\) be the essential nonperipheral simple closed curve classes and \(\ell_{X_t}(\alpha)\) their geodesic lengths. A boundary-convergence assertion in a chosen direction would mean existence of a nonzero measured lamination \(\mu\) and positive normalizing factors \(a_t\), tending to zero in that direction, such that
\[
 a_t\ell_{X_t}(\alpha)\longrightarrow i(\mu,\alpha)
 \quad\text{for every }\alpha\in\mathcal C(S),
\]
where \(i\) is geometric intersection number. This is the standard projective-length meaning of convergence, conditional on the usual surface conventions; variants with boundary may require an adjusted coordinate family.

The source asks whether stretch lines converge to individual points of this boundary and what their asymptotic behavior is. The positive and negative directions must be distinguished: \(t\to+\infty\) and \(t\to-\infty\) may require different normalizations and laminations, and a conclusion for one direction is not a conclusion for both. A unique boundary limit is stronger than existence of subsequential accumulation points. The wording also does not choose whether additional asymptotics concern curve-length growth, the stretch metric, Teichm\"uller distance, an identification of the limiting lamination or rates of convergence; no universal expansion or rate is added.

The literature triage checked 2026-08-21 reports positive-direction convergence established and important negative-direction cases known, but no complete answer for all stretch lines and directions. These are retained as an attributed partial-progress description, without an invented endpoint formula or extra lamination hypothesis. The exact-statement label and official-source verification are supplied metadata, not independently checked current status.
\subsection{Short English statement}
Ask whether each intended stretch line has a unique Thurston-boundary limit, distinguishing its positive and negative directions, and describe the relevant asymptotics. Projective length convergence gives a standard precise meaning of the endpoint question. The surface, lamination class and additional asymptotic quantities are not specified.
\subsection{Sources}
\begin{itemize}
\item Local UnsolvedMath registry, permanent ID 30000506, source code catalogue entry, ``Asymptotics of Stretch Lines in Teichm\"uller Space,'' including the unchanged original question and verification metadata.
\item \emph{Teichm\"uller Space (Classical and Quantum) (2007)}, Oberwolfach Reports, pp.\ 1537--1614, DOI 10.4171/owr/2006/26: \url{https://doi.org/10.4171/owr/2006/26}. The title uses 2007, while the DOI and proposed-year metadata use 2006.
\item Supplied reference ``Th\'eret, negative convergence'': \url{https://www.acadsci.fi/mathematica/Vol32/Theret.html}.
\item Supplied reference ``Th\'eret, stretch-line asymptotics'': \url{https://arxiv.org/abs/1804.11245}.
\item Literature triage checked 2026-08-21; record updated 2026-09-16. No external source has been visited or independently checked here.
\end{itemize}
\end{document}
